%\documentclass[10pt,respuestas,a4]{examenEPN}
\documentclass[10pt,a4]{examenEPN0}
\examen{Examen 01}
\usepackage{comandosEPN}

\usepackage{empheq}

\newcommand{\itemcp}{\stepcounter{enumi}\item[ \textcolor{red}{\theenumi.$^*$}] }

\newcommand{\iitemcp}{\stepcounter{enumii}\item[ \textcolor{red}{\theenumii)$^*$}] }


\materia{Álgebra Lineal}
%\autor{Departamento de Formación Básica}
\autor{GR \hspace{1cm}.}
%\fecha{Semestre 2026-A}
\fecha{Nombre: \hspace{9cm} Cédula:}
\logouno[1.6cm]{Logo_EPN.eps}
\logodos[2.6cm]{Logo_DFB.eps}

\newcommand{\la}{\lambda}

%Datos
\newcommand{\datosAlgebra}[2]{
\EnExamen
{
\vspace{5mm}
    \noindent
    Nombres: \hrulefill\
    Carrera: \rule{4.5cm}{0.4pt}\
    \ifthenelse{\equal{#1}{}}%
    {Paralelo: \underline{\parbox{1.8cm}{\parbox{1cm}{}}}\\[3mm]}
    {Paralelo: \underline{\parbox{1.8cm}{\centering #1}}\\[3mm]}
    C\'edula: \rule{5cm}{0.4pt}\
    Firma: \hrulefill\
    Fecha: \underline{\parbox{2.1cm}{\centering #2}}
}}


% \newcommand{\indicacionesPrueba}{%
% \EnExamen{
% \begin{indicaciones}
%     \item La prueba tiene una duraci\'on  de \textbf{120 minutos.}
%     \item No se permite la utilizaci\'on de dispositivos electr\'onicos (celular, tablets, calculadora, etc.)
%     \item Leer detenidamente los enunciados, \textbf{justificar la respuesta para acceder al puntaje total}.
%     \item La prueba es completamente individual. Los estudiantes deberán dejar sus teléfonos celulares guardados y no podrán tenerlos consigo (por ejemplo, en los bolsillos de su ropa) durante toda la evaluación. En caso de detectarse algún acto de deshonestidad académica o el incumplimiento de esta norma, no podrán seguir rindiendo la prueba y su calificación será \textbf{cero}.
% \end{indicaciones}
% \pagestyle{empty}
% }}

\newcommand{\indicacionesPrueba}{%
\EnExamen{
\begin{indicaciones}
    \item La prueba tiene una duraci\'on  de \textbf{60 minutos.}
    \item Leer detenidamente los enunciados, \textbf{justificar la respuesta} para acceder al puntaje total.
    \item No se permite la utilizaci\'on de dispositivos electr\'onicos (celular, tablets, calculadora, etc.)
    \item La prueba es completamente individual
    \item No se permite portar teléfonos celulares, relojes inteligentes ni otros dispositivos electrónicos durante el examen, aun cuando permanezcan en los bolsillos. Estos dispositivos deberán guardarse en la mochila antes del inicio de la evaluación y esta deberá mantenerse cerrada durante todo el examen.
    \item He leído las indicaciones y acepto que el incumplimiento de estas se penaliza con \textbf{cero}. Firma:...................................

    
\end{indicaciones}
\pagestyle{empty}
}}

\newcommand{\indicacionesExamen}{
\EnExamen{\begin{indicaciones}
    \item La evaluación tiene una duraci\'on  de \textbf{120 minutos.}
    \item No se permite la utilizaci\'on de dispositivos electr\'onicos (celular, tablets, calculadora, etc.)
    %\item Leer detenidamente los enunciados, justificar la respuesta para acceder al puntaje total.
    \item El examen es completamente individual. En caso de detectarse alg\'un acto de deshonestidad acad\'emica no podr\'a seguir rindiendo la prueba y su calificaci\'on ser\'a \textbf{cero}.
\end{indicaciones}
}\pagestyle{empty}
}

%% Imprimir las una hoja A4, dos veces en A5
\newcommand{\imprimirDoble}
{%
\usepackage{pgfmorepages}
\pgfmorepagesloadextralayouts
\pgfpagesuselayout{repeated 2-up}[a4paper,landscape]
\thispagestyle{empty}
}


%Este sí se puede comentar
\newenvironment{solucion}{\color{blue}\proof[Solución]}{\endproof\color{black}}




% \renewcommand{\@autor}{
% ---\\
% \phantom{.} 
% ...}
% \makeatother

%\imprimirDoble
\begin{document}
\encabezado
\datosAlgebra{GR}{18/11/2026}
%Poner el grupo y la fecha, o dejar las llaves vacías.
\indicacionesPrueba
\begin{preguntas}


%%%%%%%%%%%%%%%%%%%%%%%%%%%%%%%%%
%%%%%%%%%%%%%%%%%%%%%%%%%%%%%%%%%

\item Considere las matrices 
\[
    A=\begin{pmatrix}
        1 & \alpha & 2 \\
        2 & 3\alpha & 6\\
        3 & 4\alpha & \alpha + 2
    \end{pmatrix} \texty 
    b=\begin{pmatrix}
        1\\\alpha + 2\\\alpha + 3    
    \end{pmatrix}
\]
\begin{enumerate}
    \item Calcule $\det(A)$ y determine el valor de $\alpha\in\mathbb{R}$ para que la matriz sea invertible. \puntaje{2.0}
    \item Para $\alpha=0$, halle la solución del sistema $Ax=b$. \puntaje{2.0}
    
\end{enumerate}

%%%%%%%%%%%%%%%%%%%%%%%%%%%%%%%%%%%%%%%
\item Considere S y T subespacios vectoriales de $\R^{2\times2}$
\begin{center}
$S=\{A : a_{11}+a_{12}+a_{21}+a_{22}=0\},\qquad
T=\{A : a_{11}=a_{22},\ a_{12}=a_{21}\}.$
\end{center}
\begin{enumerate}
\item Hallar $S\cap T$ de forma explícita. \puntaje{2.0}
\item Hallar $k\in\R$ tal que $A=\begin{pmatrix}k+1&k-3\\2k-4&3k-1\end{pmatrix}\in S\cap T$. \puntaje{2.0}
\end{enumerate}

%%%%%%%%%%%%%%%%%%%%%%%%%%%%%%%%%%%%%%%%%
\item Sean \(u = (1, 2, -1)\), \(v = (2, 3, 1)\) y \(w = (1, 1, k)\) vectores de \(\mathbb{R}^3\), con \(k \in \mathbb{R}\).

\begin{enumerate}
    \item Determine \textbf{para qué valores de \(k\)} el conjunto \(\{u, v, w\}\) genera todo \(\mathbb{R}^3\). \puntaje{2.0} 
    \item Calcule \(\operatorname{span}\{u, v, w\}\) para \(k = 2\) e interprételo geométricamente. ¿Es todo \(\mathbb{R}^3\)? Justifique. \puntaje{1.0} 
    \item Para \(k = 2\), encuentre un vector en \(\operatorname{span}\{u, v, w\}\) distinto del vector nulo y de los generadores \(u, v, w\), y otro vector en \(\mathbb{R}^3\) que \textbf{no} pertenezca a \(\operatorname{span}\{u, v, w\}\). Justifique ambos casos. \puntaje{1.0} 
\end{enumerate}



%%%%%%%%%%%%%%%%%%%%%%%%%%%%%%%%%%%%%%%%%
\item Una empresa utiliza vehículos de tres tipos para un único viaje por vehículo. Cada vehículo A transporta una tonelada y requiere dos conductores; cada vehículo B transporta dos toneladas y requiere un conductor; cada vehículo C transporta tres toneladas y requiere un conductor. Ese día salieron nueve vehículos, todos a su capacidad máxima, con 17 toneladas de carga y 12 conductores en total. Determine cuántos vehículos de cada tipo salieron.

\begin{enumerate}
    \item Defina las tres variables, plantee el sistema de ecuaciones lineales que modela la situación y escriba su representación matricial $A\mathbf{x}=b$, identificando $A$, $\mathbf{x}$ y $b$. \puntaje{2.0}
    \item Resuelva el sistema e indique la cantidad de vehículos que salieron a circulación. \puntaje{2.0}
\end{enumerate}

%%%%%%%%%%%%%%%%%%%%%%%%%%%%%%%%%%%%%%%
\item (Ejercicio del deber)  Sea $V = \mathcal{P}_3(\R)$ el espacio de polinomios de grado menor o igual a tres con coeficientes reales. Determine si el siguiente conjunto es un subespacio de $V$: \puntaje{4.0}
    \[
    W = \left\{ p(x) \in \mathcal{P}_3(\R) : p(2) = 0 \quad \text{y} \quad p'(0) = 0 \right\}
    \]    




\end{preguntas}
\final
\end{document}